%% ============================================================
%% TE3001B — Fundamentación de Robótica
%% Teleoperación con Sensado de Fuerza: xArm Lite 6
%% Formato IEEE Transactions — Versión Final con Imágenes
%% Prof. Alberto Muñoz — Computational Robotics Lab
%% Tecnológico de Monterrey — 2026
%% ============================================================
\documentclass[journal,10pt,twoside]{IEEEtran}

%% ── Paquetes ─────────────────────────────────────────────────
\usepackage[utf8]{inputenc}
\usepackage[T1]{fontenc}
\usepackage[spanish,es-nodecimaldot]{babel}
\usepackage{amsmath,amssymb,bm}
\usepackage{graphicx}
\usepackage{xcolor}
\usepackage{hyperref}
\usepackage{booktabs}
\usepackage{multirow}
\usepackage{array}
\usepackage{colortbl}
\usepackage{caption}
\usepackage{subcaption}
\usepackage{listings}
\usepackage{tcolorbox}
\usepackage{enumitem}
\usepackage{float}
\usepackage{url}
\tcbuselibrary{breakable,skins}

\hypersetup{colorlinks=true,linkcolor=tec_blue,urlcolor=tec_blue,citecolor=tec_blue}

%% ── Colores ──────────────────────────────────────────────────
\definecolor{tec_blue}{RGB}{0,68,124}
\definecolor{tec_lightblue}{RGB}{0,140,200}
\definecolor{tec_navy}{RGB}{0,31,84}
\definecolor{tec_iceblue}{RGB}{220,237,250}
\definecolor{tec_paleblue}{RGB}{235,245,255}
\definecolor{tec_gray}{RGB}{100,100,100}
\definecolor{code_bg}{RGB}{13,17,23}
\definecolor{code_text}{RGB}{168,216,234}
\definecolor{code_kw}{RGB}{86,156,214}
\definecolor{code_str}{RGB}{206,145,120}
\definecolor{code_cmt}{RGB}{106,153,85}
\definecolor{box_info_border}{RGB}{0,68,124}
\definecolor{box_warn_border}{RGB}{200,140,0}

%% ── Cajas ────────────────────────────────────────────────────
\newtcolorbox{infobox}[1][Nota]{
  colback=tec_iceblue,colframe=box_info_border,
  title={\textbf{#1}},fonttitle=\bfseries\small\color{tec_navy},
  breakable,enhanced,left=4pt,right=4pt,top=3pt,bottom=3pt}
\newtcolorbox{eqbox}[1][Ecuacion]{
  colback=tec_paleblue,colframe=tec_lightblue,
  title={\textbf{#1}},fonttitle=\bfseries\small\color{tec_blue},
  breakable,enhanced,left=4pt,right=4pt}

%% ── Código ───────────────────────────────────────────────────
\lstdefinestyle{python}{
  language=Python,backgroundcolor=\color{code_bg},
  basicstyle=\ttfamily\scriptsize\color{code_text},
  keywordstyle=\color{code_kw}\bfseries,
  stringstyle=\color{code_str},commentstyle=\color{code_cmt}\itshape,
  numbers=left,numberstyle=\tiny\color{tec_gray},numbersep=5pt,
  frame=single,framerule=0.5pt,rulecolor=\color{tec_blue},
  tabsize=4,showstringspaces=false,breaklines=true,
  literate={á}{{\'{a}}}1{é}{{\'{e}}}1{í}{{\'{i}}}1
           {ó}{{\'{o}}}1{ú}{{\'{u}}}1{ñ}{{\~{n}}}1}
\lstdefinestyle{bash}{
  language=bash,backgroundcolor=\color{code_bg},
  basicstyle=\ttfamily\scriptsize\color{code_text},
  keywordstyle=\color{code_kw}\bfseries,commentstyle=\color{code_cmt},
  frame=single,framerule=0.5pt,rulecolor=\color{tec_blue},
  numbers=none,breaklines=true}
\lstdefinestyle{json}{
  backgroundcolor=\color{code_bg},
  basicstyle=\ttfamily\scriptsize\color{code_text},
  frame=single,framerule=0.5pt,rulecolor=\color{tec_lightblue},
  numbers=none,breaklines=true}
\lstset{style=python}

%% ── Macros ───────────────────────────────────────────────────
\newcommand{\xarm}{xArm Lite 6}
\newcommand{\ros}{ROS2}

%% ── Metadatos ────────────────────────────────────────────────
\begin{document}

\title{\color{tec_navy}Teleoperación con Sensado de Fuerza Física\\
  entre Dos \xarm{}:\\
  Diseño, Implementación y Evaluación Experimental}

\author{%
  \IEEEauthorblockN{%
    Fabricio Banda\IEEEauthorrefmark{1},\;
    Manuel Ferro\IEEEauthorrefmark{1},
    Alexandro Cardenas\IEEEauthorrefmark{1},
    Fernando Proal\IEEEauthorrefmark{1},
    Zacbe Ortega\IEEEauthorrefmark{1}}\\
  \IEEEauthorblockA{%
    \IEEEauthorrefmark{1}Computational Robotics Lab,
    Tecnológico de Monterrey, Monterrey, N.L., México\\
    \texttt{\{A00839916,A00836182,A01385387,A01385446,A00574358\}@tec.mx}}}
\maketitle

%% ── Abstract ─────────────────────────────────────────────────
\begin{abstract}
Se presenta el diseño e implementación de un sistema de teleoperación
maestro-esclavo utilizando dos manipuladores \xarm{}.
El brazo maestro es controlado mediante sensores de presión
artesanales de esponja conductora de grafito (6~direcciones),
permitiendo al operador comandar velocidades cartesianas en $X$,
$Y$ y $Z$ manteniendo la orientación del efector fija.
El brazo esclavo replica el movimiento del maestro en tiempo real
y porta cinco sensores adicionales de fuerza que, al superar un
umbral de $|\bm{F}|>3{,}0$\,N, activan un bloqueo sincronizado
que detiene ambos brazos hasta que el contacto desaparece.
La arquitectura se implementó sobre \ros{} (Jazzy), calculando el
Jacobiano geométrico dinámicamente desde el árbol TF y aplicando
pseudoinversa amortiguada (DLS, $\lambda=0{,}01$).
Se realizaron tres pares de experimentos en simulación y en los
robots físicos, obteniendo errores de seguimiento articular de
$0{,}33$--$1{,}10$°\,RMS y un retardo de seguimiento de
$\approx$\,100--150\,ms en condiciones reales.
Video de demostración disponible en: \url{https://youtu.be/9yoCLF5ZJVQ}
\end{abstract}

\begin{IEEEkeywords}
  teleoperación, sensado de fuerza, \xarm{}, Jacobiano geométrico,
  ROS2, ESP32, sensores de grafito, pseudoinversa amortiguada,
  detección de colisión, maestro-esclavo
\end{IEEEkeywords}

%% ============================================================
\section{Introducción}
%% ============================================================

\IEEEPARstart{L}{a} teleoperación robótica con retroalimentación de
fuerza es fundamental para extender las capacidades humanas en
entornos remotos o peligrosos~\cite{murray1994}. Aplicaciones como
cirugía mínimamente invasiva~\cite{hogan1985}, exploración espacial
y manipulación de materiales peligrosos requieren que el operador
``sienta'' el entorno remoto a través del robot.

El \xarm{}, desarrollado por UFACTORY~\cite{xarm_github}, es un
manipulador de 6~DOF para investigación. Su SDK Python permite
control de posición, velocidad articular y modo servo en tiempo real.

\subsection{Motivación}
Se diseñaron sensores de presión artesanales de \emph{esponja
conductora de grafito} como alternativa económica a los sensores
F/T comerciales (que pueden costar miles de dólares), explorando
la viabilidad en robótica educativa.

\subsection{Contribuciones}
\begin{itemize}[leftmargin=1.2em,itemsep=1pt]
  \item Sensores de grafito para detección de fuerza en 5 y 6 direcciones.
  \item Nodo \ros{} para teleoperación cartesiana con Jacobiano
        calculado dinámicamente desde TF.
  \item Detección de colisión en tiempo real con bloqueo sincronizado
        e histéresis configurable.
  \item Visualización en tiempo real de $F_x$, $F_y$, $F_z$ y $|\bm{F}|$.
  \item Integración completa ESP32 + serial JSON + \ros{} Jazzy.
\end{itemize}

%% ============================================================
\section{Marco Teórico}
%% ============================================================

\subsection{Jacobiano Geométrico}

El \xarm{} tiene $n=6$ DOF con articulaciones rotacionales.
El Jacobiano $\bm{J}(\bm{q})\in\mathbb{R}^{6\times6}$ mapea
velocidades articulares a cartesianas~\cite{siciliano2009}:
\begin{equation}
  \dot{\bm{x}} = \bm{J}(\bm{q})\,\dot{\bm{q}}
  \label{eq:jacobian}
\end{equation}
Para cada articulación rotacional~$i$:
\begin{equation}
  \bm{J}_i =
  \begin{pmatrix}
    \bm{z}_i \times (\bm{p}_{ee} - \bm{p}_i) \\
    \bm{z}_i
  \end{pmatrix}
  \label{eq:jcol}
\end{equation}

\subsection{Pseudoinversa Amortiguada (DLS)}

Para la cinemática inversa diferencial se usa DLS~\cite{spong2006}:
\begin{equation}
  \bm{J}^{+} = \bm{J}^\top
    \!\left(\bm{J}\bm{J}^\top + \lambda^2\bm{I}\right)^{-1},
  \quad \lambda=0{,}01
  \label{eq:dls}
\end{equation}
El incremento articular por ciclo es $\Delta\bm{q}=\bm{J}^{+}\bm{\xi}$,
donde $\bm{\xi}=[v_x,v_y,v_z,0,0,0]^\top$ fuerza traslación pura.

\subsection{Elección del Modo de Control: SERVO vs. Par Calculado}
\label{sec:ctc}

El control de \textbf{par calculado} (CTC) es el esquema de control
óptimo para un manipulador rígido, pues linealiza y desacopla la dinámica
del robot mediante la ley:
\begin{equation}
  \bm{\tau} = \bm{M}(\bm{q})\bigl(K_p\bm{e}+K_v\dot{\bm{e}}+\ddot{\bm{q}}_d\bigr)
            + \bm{C}(\bm{q},\dot{\bm{q}})\dot{\bm{q}} + \bm{g}(\bm{q})
  \label{eq:ctc}
\end{equation}
produciendo la dinámica de error lineal
$\ddot{\bm{e}} + K_v\dot{\bm{e}} + K_p\bm{e} = \bm{0}$~\cite{siciliano2009}.

Sin embargo, implementar CTC sobre el \xarm{} desde la SDK Python
presenta dos obstáculos técnicos insuperables en este proyecto:

\begin{enumerate}[leftmargin=1.4em,itemsep=2pt]
  \item \textbf{Acceso limitado al modo torque.}
        La SDK del \xarm{} expone el modo de control de torque
        (\texttt{set\_mode(4)}) pero la latencia de la interfaz TCP
        interna es de $\approx$\,10--50\,ms por comando, lo que
        hace imposible cerrar el lazo de par a la frecuencia mínima
        requerida por CTC ($>$\,500\,Hz).
  \item \textbf{Modelo dinámico no disponible.}
        Las matrices $\bm{M}(\bm{q})$, $\bm{C}(\bm{q},\dot{\bm{q}})$
        y $\bm{g}(\bm{q})$ del \xarm{} no son públicas.
        Identificarlas experimentalmente requeriría un proceso de
        caracterización dinámica fuera del alcance de este proyecto.
\end{enumerate}

Por estas razones se optó por el \textbf{modo SERVO}
(\texttt{set\_mode(1)} + \texttt{set\_servo\_angle}), que delega el
control de posición articular al controlador interno propietario del
robot, operando a su frecuencia nominal de $\approx$\,250\,Hz de forma
transparente. El lazo externo implementado en Python opera a 10--20\,Hz
y envía únicamente referencias de posición $\bm{q}^{\text{ref}}$,
lo que equivale funcionalmente a un \textbf{control de posición en
espacio articular}. Esta elección sacrifica el desacoplamiento dinámico
de CTC, pero garantiza estabilidad y aprovecha el controlador
ya sintonizado de fábrica por UFACTORY.

\begin{infobox}[Justificacion del modo SERVO]
  \small
  El modo SERVO del \xarm{} implementa internamente un controlador
  de posición/velocidad de alta frecuencia. Al enviarle referencias
  a 10--20\,Hz desde Python, el sistema opera en cascada:
  lazo externo (ROS2, 10--20\,Hz) $\to$ lazo interno (firmware,
  $\approx$\,250\,Hz). Este esquema es equivalente al usado en
  sistemas comerciales de teleoperación de baja latencia como
  los reportados en~\cite{niemeyer2008}.
\end{infobox}

\subsection{Principio Piezorresistivo}

La esponja de grafito reduce su resistencia al comprimirse.
La presión normalizada leída por el ADC del ESP32 es:
\begin{equation}
  p_i = \frac{\text{BASELINE}_i - \text{ADC}_i}
             {\text{BASELINE}_i - \text{PRESSED}_i},
  \quad p_i\in[0,1]
  \label{eq:pressure}
\end{equation}
con zona muerta $p_i<0{,}08\Rightarrow p_i=0$.

\subsection{Estimación de Fuerza por Diferencial}
\begin{align}
  F_x &= (p_{\text{R}}-p_{\text{L}})\cdot15{,}0\;\text{N}  \label{eq:fx}\\
  F_y &= (p_{\text{F}}-p_{\text{B}})\cdot15{,}0\;\text{N}  \label{eq:fy}\\
  F_z &= p_{\text{Bot}}\cdot15{,}0\;\text{N}               \label{eq:fz}\\
  |\bm{F}| &= \sqrt{F_x^2+F_y^2+F_z^2}                    \label{eq:fm}
\end{align}

%% ============================================================
\section{Arquitectura del Sistema}
\label{sec:arch}
%% ============================================================

\subsection{Visión General}

El sistema integra: (1)~brazo maestro con sensores de presión de
6~direcciones, (2)~brazo esclavo con sensores de fuerza de
5~direcciones, (3)~dos ESP32, y (4)~\ros{} como middleware.
La Fig.~\ref{fig:bloques} muestra la arquitectura completa.

\begin{figure}[H]
  \centering
  \includegraphics[width=\columnwidth]{figs/diagram_bloques.jpg}
  \caption{Diagrama de bloques del sistema de teleoperación
    maestro-esclavo. Los tópicos \ros{} en amarillo conectan los
    nodos de la capa central. Las bandas de color indican el
    subsistema al que pertenece cada bloque.}
  \label{fig:bloques}
\end{figure}

\subsection{Topología de Nodos \ros{}}

\begin{table}[H]
  \centering\footnotesize\caption{Nodos \ros{} del sistema}
  \label{tab:nodes}
  \begin{tabular}{p{2.1cm}p{2.4cm}p{1.4cm}p{1.0cm}}
    \toprule
    \textbf{Nodo}&\textbf{Función}&\textbf{Publica}&\textbf{Suscribe}\\
    \midrule
    \rowcolor{tec_iceblue}
    \texttt{master\_input}&Control Jacobiano&\texttt{/L\_traj}&\texttt{/joint\_states}\\
    \texttt{dual\_mirror}&Espejo + colisión&\texttt{/R\_traj /master\_hold}&\texttt{/joint\_states /slave/force}\\
    \rowcolor{tec_iceblue}
    \texttt{force\_reader5}&Lectura serial&\texttt{/slave/force}&---\\
    \texttt{force\_display5}&Visualización&---&\texttt{/slave/force}\\
    \bottomrule
  \end{tabular}
\end{table}

\subsection{Flujo de Datos}

En cada ciclo de 50\,ms (20\,Hz):
\begin{enumerate}[leftmargin=1.4em,itemsep=1pt]
  \item El ESP32 maestro envía \texttt{\{vx,vy,vz\}} por serial.
  \item \texttt{master\_input} calcula $\bm{J}$ desde TF, computa
        $\Delta\bm{q}$ y publica la trayectoria al maestro.
  \item \texttt{dual\_mirror} replica las posiciones al esclavo.
  \item El ESP32 esclavo envía \texttt{\{fx,fy,fz,fm\}} por serial.
  \item Si $|\bm{F}|>3{,}0$\,N, publica
        \texttt{/master\_hold=True} y ambos brazos se detienen.
\end{enumerate}

%% ============================================================
\section{Diseño de Sensores de Presión}
\label{sec:sensors}
%% ============================================================

\subsection{Construcción Física}

Cada sensor artesanal consiste en una esponja de
$\approx$1\,cm$^3$ impregnada con grafito en sus caras,
dos cables conductores y cinta adhesiva de protección.
La Fig.~\ref{fig:sensor_individual} muestra un sensor individual.

\begin{figure}[H]
  \centering
  \includegraphics[width=0.7\columnwidth]{figs/foto_sensor_individual.jpg}
  \caption{Sensor de presión artesanal de esponja-grafito.
    El bloque blanco es la esponja encapsulada con cinta;
    los cables morado y blanco conectan al divisor resistivo
    del ESP32.}
  \label{fig:sensor_individual}
\end{figure}

\subsection{Arreglo de 5 Sensores con ESP32}

La Fig.~\ref{fig:sensores_esp32} muestra los cinco sensores del
esclavo conectados a la protoboard y al ESP32, con cada sensor
etiquetado según su dirección (R, L, F, B, Bottom).

\begin{figure}[H]
  \centering
  \includegraphics[width=\columnwidth]{figs/foto_sensores_esp32.jpg}
  \caption{Arreglo completo de 5 sensores de fuerza conectados a
    la ESP32 (lado esclavo). Cada sensor va a un pin ADC1
    independiente a través de un divisor resistivo de 10\,k$\Omega$.
    Se usa la misma configuración (con 6 sensores) para el maestro.}
  \label{fig:sensores_esp32}
\end{figure}

\subsection{Montaje en el Efector Final}

Las Figs.~\ref{fig:montaje1} y~\ref{fig:montaje2} muestran los
sensores montados sobre el efector del brazo esclavo, orientados
en las cinco direcciones cartesianas de detección.

\begin{figure}[H]
  \centering
  \begin{subfigure}[b]{0.48\columnwidth}
    \includegraphics[width=\textwidth]{figs/foto_slave_montado1.jpg}
    \caption{Vista frontal del brazo esclavo con el arreglo
      de sensores y el ESP32 montados.}
    \label{fig:montaje1}
  \end{subfigure}\hfill
  \begin{subfigure}[b]{0.48\columnwidth}
    \includegraphics[width=\textwidth]{figs/foto_slave_montado2.jpg}
    \caption{Vista del efector con sensores en las
      direcciones Right, Left, Forward y Back.}
    \label{fig:montaje2}
  \end{subfigure}
  \caption{Montaje de los sensores de fuerza en el efector del
    \xarm{} esclavo. El ESP32 va sujeto al enlace L2 del brazo
    para minimizar el peso en el efector.}
  \label{fig:montaje}
\end{figure}

\subsection{Firmware ESP32 y Calibración}

\begin{lstlisting}[style=json,caption={JSON publicado por ESP32 esclavo a 20 Hz}]
{"r":[2100,890,1390,2750,1180],
 "p":[0.081,0.000,0.000,0.234,0.000],
 "fx":1.23,"fy":-3.51,"fz":0.00,"fm":3.72}
\end{lstlisting}

\begin{lstlisting}[caption={Normalización y suavizado en firmware}]
float range = (float)(BASELINE[i] - PRESSED[i]);
float val   = (float)(BASELINE[i] - raw[i]) / range;
val = constrain(val, 0.0f, 1.0f);
if (val < DEADZONE) val = 0.0f;   // zona muerta 8%
smoothed[i] = ALPHA*val+(1.0f-ALPHA)*smoothed[i];
\end{lstlisting}

La Fig.~\ref{fig:terminal_arduino} muestra la salida de la terminal
de Arduino en modo calibración cuando el operador presiona el sensor
de dirección \emph{Forward} (\texttt{Fy+}).

\begin{figure}[H]
  \centering
  \includegraphics[width=\columnwidth]{figs/foto_terminal_arduino.jpg}
  \caption{Salida de la terminal de Arduino (modo calibración)
    mostrando la respuesta JSON al presionar el sensor \emph{Forward}.
    Se observa cómo \texttt{p[2]} aumenta de 0 a $\approx$\,0{,}83
    mientras el resto permanece en 0, y \texttt{fy} sube hasta
    $\approx$\,12{,}4\,N. El campo \texttt{fm} refleja la magnitud
    total de la fuerza.}
  \label{fig:terminal_arduino}
\end{figure}

%% ============================================================
\section{Implementación de Software}
\label{sec:impl}
%% ============================================================

\subsection{Nodo \texttt{master\_input.py}}

Recibe presiones del ESP32, calcula el Jacobiano desde TF y
envía trayectorias articulares.

\begin{lstlisting}[caption={Jacobiano + DLS en master\_input.py}]
def _compute_jacobian_6x6(self):
    p_ee, _, _ = self._lookup(EE_FRAME, BASE_FRAME)
    J = np.zeros((6, 6))
    for i, frame in enumerate(self.joint_frames):
        p_i, R_i, _ = self._lookup(frame, BASE_FRAME)
        z_i = R_i[:, 2]
        J[0:3,i] = np.cross(z_i, p_ee - p_i)
        J[3:6,i] = z_i
    return J

# DLS (lambda=0.01)
JJT    = J @ J.T + 0.01**2 * np.eye(6)
J_pinv = J.T @ np.linalg.inv(JJT)
dq     = J_pinv @ np.array([vx,vy,vz,0,0,0])
q_new  = np.clip(current_joints + dq, q_min, q_max)
\end{lstlisting}

\subsection{Nodo \texttt{dual\_mirror.py} — Detección de Colisión}

\begin{lstlisting}[caption={Lógica de contacto con histéresis}]
def force_cb(self, msg):
    self.fm = sqrt(fx**2 + fy**2 + fz**2)
    if (not self.hold) and self.fm >= 3.0:  # activar
        self.hold = True
        hold_msg = Bool(); hold_msg.data = True
        self.master_hold_pub.publish(hold_msg)
    elif self.hold and self.fm <= 2.0:      # liberar
        self.hold = False
        hold_msg = Bool(); hold_msg.data = False
        self.master_hold_pub.publish(hold_msg)
\end{lstlisting}

\subsection{Contact Hold como Retroalimentación Háptica Simplificada}
\label{sec:haptic}

En un sistema háptico completo, la retroalimentación al operador se
implementa aplicando torques articulares en el maestro proporcionales
a la fuerza medida en el esclavo~\cite{hogan1985}:
\begin{equation}
  \bm{\tau}_{m,\text{fb}}
    = \bm{J}_m^\top \cdot \beta \cdot \hat{\bm{F}}_{\text{ext}}
  \label{eq:haptic}
\end{equation}
donde $\beta$ es el factor de escala de retroalimentación y
$\hat{\bm{F}}_{\text{ext}}$ la fuerza estimada en el esclavo.
Esta implementación requiere que el maestro opere en modo torque
con acceso directo a $\bm{\tau}$, lo que no es factible con la
SDK del \xarm{} a las frecuencias necesarias (ver
Sección~\ref{sec:ctc}).

Como alternativa, el sistema implementa \textbf{retroalimentación
háptica posicional}: al detectar contacto ($|\bm{F}|>3{,}0$\,N),
el maestro se congela en su posición actual publicando
\texttt{/master\_hold=True}.

Desde la perspectiva del control de impedancia~\cite{hogan1985},
este comportamiento es equivalente a aplicar una rigidez virtual
infinita ($K_d\to\infty$) en el maestro al momento del contacto:
el operador ``siente'' el obstáculo porque su brazo deja de responder
a sus comandos hasta que la fuerza desaparece, replicando la
sensación de haber ``chocado contra una pared'' sin necesidad de
torques activos en el actuador. Esta técnica, aunque más simple que
la retroalimentación de fuerza continua, es suficiente para el
objetivo principal del proyecto: evitar que el operador empuje al
esclavo contra un obstáculo sin darse cuenta.

La histéresis entre el umbral de activación (3{,}0\,N) y el de
liberación (2{,}0\,N) previene oscilaciones rápidas de bloqueo/
desbloqueo cuando la fuerza fluctúa cerca del umbral, mejorando
la estabilidad perceptual del operador.

\begin{infobox}[Comparativa con retroalimentacion de fuerza activa]
  \small
  \begin{tabular}{p{2.0cm}p{2.0cm}p{2.0cm}}
    \toprule
    \textbf{Propiedad} & \textbf{Ec.~\eqref{eq:haptic} (ideal)} & \textbf{Contact hold (impl.)}\\
    \midrule
    Retroalimentación & Continua, proporcional & Binaria (freeze/libre)\\
    Frecuencia & $>500$\,Hz & 10--20\,Hz\\
    Modo xArm & Torque (\texttt{mode=4}) & Servo (\texttt{mode=1})\\
    Estabilidad & Requiere análisis $\beta_{max}$ & Inherente\\
    Costo impl. & Alto & Bajo\\
    \bottomrule
  \end{tabular}
\end{infobox}

\begin{lstlisting}[style=bash,caption={Secuencia de lanzamiento}]
# T1: robots reales (o fake para simulacion)
ros2 launch xarm_moveit_config \
  dual_lite6_moveit_realmove.launch.py \
  robot_ip_1:=192.168.1.179 robot_ip_2:=192.168.1.154
# T2: espejo + colision
python3 dual_mirror.py --force-stop-threshold 3.0
# T3: control maestro
python3 master_input.py --port /dev/ttyUSB1
# T4: sensores fuerza esclavo
python3 force_reader5.py --port /dev/ttyUSB0
# T5: visualizacion (opcional)
python3 force_display5.py
\end{lstlisting}

%% ============================================================
\section{Visualización de Fuerza en Tiempo Real}
\label{sec:display}
%% ============================================================

El nodo \texttt{force\_display5.py} proporciona tres paneles
actualizados a 20\,Hz: (1)~dirección de fuerza en el plano XY,
(2)~barras por canal, y (3)~historial temporal de 10\,s.

\begin{figure}[H]
  \centering
  \includegraphics[width=\columnwidth]{figs/foto_display_libre.jpg}
  \caption{Estado \emph{sin contacto}: todas las barras en cero,
    $|\bm{F}|=0{,}0$\,N, sin flecha en el panel XY. El umbral de
    detección (línea roja punteada a 2{,}5\,N) no ha sido superado.}
  \label{fig:display_libre}
\end{figure}

\begin{figure}[H]
  \centering
  \includegraphics[width=\columnwidth]{figs/foto_display_bajo_umbral.jpg}
  \caption{Fuerza presente ($|\bm{F}|\approx1{,}9$\,N) pero
    \textbf{por debajo del umbral}: el sistema reporta
    ``Sin contacto'' aunque $F_y\approx-1{,}9$\,N es visible en la
    barra \texttt{Fy-}. Esto confirma que la detección no es On/Off
    sino comparada con un umbral configurable.}
  \label{fig:display_bajo_umbral}
\end{figure}

\begin{figure}[H]
  \centering
  \includegraphics[width=\columnwidth]{figs/foto_display_contacto.jpg}
  \caption{Contacto detectado en dirección \emph{Forward}:
    $F_y=+15{,}0$\,N, $|\bm{F}|=15{,}0$\,N, superando el umbral
    de 2{,}5\,N. El panel muestra ``!! CONTACTO !!'' en rojo,
    la flecha apunta al norte (FWD) y la barra \texttt{Fy+}
    está saturada en rojo.}
  \label{fig:display_contacto}
\end{figure}

\begin{figure}[H]
  \centering
  \includegraphics[width=\columnwidth]{figs/foto_display_dos_fuerzas.jpg}
  \caption{Contacto simultáneo en dos direcciones:
    $F_x=-15{,}0$\,N (izquierda) y $F_y=-12{,}0$\,N (atrás),
    con $|\bm{F}|=19{,}2$\,N. La flecha del panel XY apunta en la
    dirección resultante (diagonal inferior-izquierda), y las barras
    \texttt{Fx-} y \texttt{Fy-} están activas simultáneamente.
    Esto valida la estimación vectorial de la fuerza.}
  \label{fig:display_dos_fuerzas}
\end{figure}

%% ============================================================
\section{Setup Experimental}
\label{sec:setup}
%% ============================================================

\begin{figure}[H]
  \centering
  \includegraphics[width=\columnwidth]{figs/foto_setup_completo.jpg}
  \caption{Setup experimental completo en el Computational
    Robotics Lab: el brazo maestro (derecha) y el brazo esclavo
    (izquerda) sobre la misma mesa de trabajo, con la laptop
    central ejecutando todos los nodos \ros{} y el visualizador.}
  \label{fig:setup}
\end{figure}

\begin{infobox}[Video de demostracion y repositorio]
  \small
  \textbf{Video} — teleoperación en tiempo real con ambos
  \xarm{} y evento de \textit{contact hold}:\\
  \url{https://youtu.be/9yoCLF5ZJVQ}\\[3pt]
  \textbf{Repositorio GitHub} — código, firmware ESP32, datos de
  experimentos y reporte LaTeX{}:\\
  \url{https://github.com/kurttco/TeleOp-Haptica.git}\\
\end{infobox}
  \centering\footnotesize\caption{Especificaciones del hardware}
  \label{tab:hw}
  \setlength{\tabcolsep}{4pt}
  \begin{tabular}{lll}
    \toprule
    \textbf{Componente}&\textbf{PC Maestro}&\textbf{PC Esclavo}\\
    \midrule
    \rowcolor{tec_iceblue}
    OS & Ubuntu 24 + ROS2 Jazzy & Ubuntu 24 + ROS2 Jazzy\\
    Python & 3.10+ & 3.10+\\
    \rowcolor{tec_iceblue}
    IP xArm & 192.168.1.179 & 192.168.1.154\\
    ESP32 port & /dev/ttyUSB1 & /dev/ttyUSB0\\
    \rowcolor{tec_iceblue}
    Sensores & 6 direcciones & 5 direcciones\\
    \bottomrule
  \end{tabular}
\end{table}

%% ============================================================
\section{Resultados Experimentales}
\label{sec:exp}
%% ============================================================

\subsection{Descripción de los Experimentos}

Se realizaron tres pares de experimentos, ejecutando cada
trayectoria tanto en simulación (\texttt{dual\_lite6\_moveit\_fake})
como con los robots físicos. El operador intentó reproducir la misma
trayectoria en ambos entornos para una comparación válida.
La Tabla~\ref{tab:exp_overview} resume las condiciones.

\begin{table}[H]
  \centering\footnotesize\caption{Resumen de los seis experimentos}
  \label{tab:exp_overview}
  \setlength{\tabcolsep}{4pt}
  \begin{tabular}{llccc}
    \toprule
    \textbf{Par}&\textbf{Cond.}&\textbf{Dur.}&\textbf{Muestras}&\textbf{Holds}\\
    \midrule
    \rowcolor{tec_iceblue}
    \multirow{2}{*}{1}&Sim & 342\,s & 740  & 1\\
                       &Real& 287\,s & 2870 & 4\\
    \rowcolor{tec_paleblue}
    \multirow{2}{*}{2}&Sim &  92\,s & 769  & 0\\
                       &Real& 118\,s & 1180 & 3\\
    \rowcolor{tec_iceblue}
    \multirow{2}{*}{3}&Sim &  40\,s & 304  & 0\\
                       &Real&  62\,s &  620 & 2\\
    \bottomrule
  \end{tabular}
\end{table}

\subsection{Seguimiento Articular — Simulación vs. Real}

La Fig.~\ref{fig:tracking} muestra las trayectorias maestro y
esclavo superpuestas para los seis experimentos.

\begin{figure*}[!t]
  \centering
  \includegraphics[width=\textwidth]{figs/fig1_tracking.jpg}
  \caption{Seguimiento articular maestro-esclavo para los tres
    pares de experimentos (columnas izquierdas\,=\,simulación en
    azul; columnas derechas\,=\,robot real en rojo). Las bandas
    amarillas indican los instantes de \textit{contact hold} en los
    experimentos reales. En la simulación, maestro y esclavo
    siguen trayectorias casi idénticas con mínima separación visual;
    en el robot real se aprecia un ligero desfase temporal de
    $\approx$\,150\,ms y una oscilación de pequeña amplitud.}
  \label{fig:tracking}
\end{figure*}

\subsection{Error de Seguimiento}

La Fig.~\ref{fig:error} muestra el error articular
$e_i(t)=q_i^m(t)-q_i^s(t)$ para los seis joints.

\begin{figure*}[!t]
  \centering
  \includegraphics[width=\textwidth]{figs/fig2_error.jpg}
  \caption{Error de seguimiento $e_i(t)=q_i^{\text{maestro}}-q_i^{\text{esclavo}}$
    en grados. Azul\,=\,simulación; naranja\,=\,robot real.
    Las zonas amarillas indican los eventos de
    \textit{contact hold}. El pico de 31° en J3 de la simulación
    (Par\,1) corresponde al artefacto de snap post-hold del
    simulador; en el robot real el controlador interno interpola
    esa transición distribuyéndola en $\approx$\,3--5 muestras.}
  \label{fig:error}
\end{figure*}

\subsubsection{Comportamiento en Simulación}

El esclavo replica al maestro con alta fidelidad debido a la
latencia de loopback ($<$\,1\,ms) y al modelo dinámico idéntico.
El error $e_i(t)$ permanece en $\pm 2$--$5$° durante el movimiento
continuo. La única excepción es el pico de 31° en J3 del Par~1
al liberar el contact hold: el nodo \texttt{dual\_mirror} envía
la posición del maestro sin interpolación y el simulador la acepta
instantáneamente.

\subsubsection{Comportamiento en el Robot Real}

Dos diferencias principales respecto a la simulación:

\begin{enumerate}[leftmargin=1.4em,itemsep=2pt]
  \item \textbf{Retardo de 100--150\,ms.} La cadena de latencias
        incluye: lectura de \texttt{/joint\_states} ($\approx$\,10\,ms),
        publicación TCP/Ethernet ($\approx$\,15\,ms) y procesamiento
        interno del controlador \xarm{} ($\approx$\,100\,ms en
        modo SERVO).
  \item \textbf{Oscilaciones($\pm$\,0{,}5--2°).}
        Causada por: (a)~ruido del encoder ($\approx$\,0{,}1°),
        (b)~oscilaciones del PID interno ante referencias rápidas,
        y (c)~eventos de slip mecánico con decaimiento exponencial.
\end{enumerate}

\subsubsection{Articulación J4}

En los pares~1 y~3, J4 casi no se mueve ($<$\,0{,}5°). En
simulación el error es prácticamente cero; en el real
$\approx$\,0{,}1° por vibraciones transmitidas de otras
articulaciones. En el Par~2 (rango $\approx$\,41°), los errores
son comparables entre ambas condiciones.

\subsection{Métricas de Desempeño}

\begin{table}[H]
  \centering\footnotesize
  \caption{Error RMS [°] por articulación y condición}
  \label{tab:metrics}
  \setlength{\tabcolsep}{3pt}
  \begin{tabular}{llccccccc}
    \toprule
    \multirow{2}{*}{\textbf{Par}}&\multirow{2}{*}{\textbf{Cond.}}&
    \multicolumn{6}{c}{\textbf{RMS por articulación [°]}}&\textbf{Global}\\
    \cmidrule(lr){3-8}
    &&J1&J2&J3&J4&J5&J6&\\
    \midrule
    \rowcolor{tec_iceblue}
    \multirow{2}{*}{1}&Sim &0.37&1.27&1.77&0.00&1.17&0.37&1.03\\
    &Real&0.15&0.41&0.54&0.11&0.39&0.15&0.33\\
    \rowcolor{tec_paleblue}
    \multirow{2}{*}{2}&Sim &0.91&0.89&0.75&0.50&0.68&0.59&0.73\\
    &Real&0.56&0.46&0.39&0.27&0.36&0.44&0.42\\
    \rowcolor{tec_iceblue}
    \multirow{2}{*}{3}&Sim &0.53&1.66&0.45&0.02&1.32&0.52&0.93\\
    &Real&0.25&1.97&0.25&0.10&1.80&0.26&1.10\\
    \bottomrule
  \end{tabular}
\end{table}

La Figs.~\ref{fig:metrics} y~\ref{fig:summary} presentan
la comparativa gráfica.

\begin{figure*}[!t]
  \centering
  \includegraphics[width=\textwidth]{figs/fig3_metrics.jpg}
  \caption{Comparativa de error RMS (fila superior) y error
    máximo (fila inferior) por articulación para los tres pares.
    Azul\,=\,simulación; rojo\,=\,robot real. Los valores $\Delta$
    en verde indican que el robot real fue mejor; en rojo, que fue
    peor. Nótese que en el Par\,1 el simulador tiene errores
    significativamente mayores en J2, J3 y J5 por el artefacto
    del snap post-hold.}
  \label{fig:metrics}
\end{figure*}

\begin{figure}[H]
  \centering
  \includegraphics[width=\columnwidth]{figs/fig4_summary.jpg}
  \caption{\textit{Izquierda:} RMS global por par.
    \textit{Derecha:} mapa de calor del RMS por articulación;
    tonos más oscuros indican mayor error. J2, J3 y J5 son
    consistentemente las articulaciones con mayor error por
    su mayor rango y velocidad; J1, J4 y J6 permanecen bajos.}
  \label{fig:summary}
\end{figure}

\subsection{Análisis: ¿Por qué RMS real $<$ RMS simulación?}

\begin{infobox}[Explicación del resultado contraintuitivo]
  \small
  En los pares~1 y~2, el RMS global del robot real
  (0{,}33--0{,}42°) es \textbf{menor} que el de la simulación
  (0{,}73--1{,}03°). Esto se debe a dos causas:

  \textbf{(1) Snap post-hold en simulación.}
  Al liberar un contact hold, el simulador ejecuta un salto
  instantáneo de hasta 2{,}6\,rad en 0{,}1\,s (26\,rad/s,
  imposible físicamente), que infla el RMS de toda la sesión.

  \textbf{(2) Mayor duración de las sesiones simuladas.}
  La sesión simulada del Par\,1 dura 342\,s vs.\ 287\,s del real,
  acumulando más tramos de movimiento rápido donde el error es mayor.

  En el \textbf{Par\,3} el robot real sí supera al simulado
  (1{,}10° vs.\ 0{,}93°), que es el comportamiento canónico.
  Para un análisis justo, el simulador debería implementar
  interpolación de trayectoria al salir del hold.
\end{infobox}

\subsection{Tendencias Globales}

El mapa de calor (Fig.~\ref{fig:summary}) revela dos patrones
transversales:
\begin{itemize}[leftmargin=1.2em,itemsep=2pt]
  \item \textbf{J2, J3, J5} tienen los mayores errores en ambas
        condiciones: son las articulaciones de mayor rango y
        velocidad, por lo que el retardo de 100--150\,ms genera
        errores proporcionales.
  \item \textbf{J1, J4, J6} permanecen bajos ($<$\,1°):
        J1 (base) se mueve lento; J4 tiene rango limitado en
        estas trayectorias; J6 (rotación efector) se fija a
        cero por diseño.
\end{itemize}

%% ============================================================
\section{Discusión}
\label{sec:disc}
%% ============================================================

\subsection{Limitaciones del Sistema}
\begin{itemize}[leftmargin=1.2em,itemsep=2pt]
  \item \textbf{Deriva de sensores:} el BASELINE de la esponja
        se desplaza con temperatura y envejecimiento, requiriendo
        recalibración periódica.
  \item \textbf{Frecuencia de control (10--20\,Hz):} baja respecto
        a sistemas comerciales ($>$\,1\,kHz); limita la velocidad
        máxima segura.
  \item \textbf{Snap post-hold en simulador:} sin interpolación
        al salir del hold, el error simulado se infla artificialmente.
  \item \textbf{Detección binaria de contacto:} no cuantifica
        fuerzas calibradas en Newtons absolutos.
\end{itemize}

\subsection{Trabajo Futuro}
\begin{itemize}[leftmargin=1.2em,itemsep=2pt]
  \item Integrar un sensor F/T comercial para validar la
        estimación artesanal.
  \item Implementar control de impedancia virtual~\cite{hogan1985}
        en lugar de parada abrupta.
  \item Añadir retroalimentación háptica activa:
        $\bm{\tau}_{m,fb}=\bm{J}_m^\top\cdot\beta\cdot\hat{\bm{F}}_{ext}$.
  \item Implementar interpolación de trayectoria al salir del
        contact hold en el simulador.
\end{itemize}

%% ============================================================
\section{Conclusiones}
\label{sec:conc}
%% ============================================================

\IEEEPARstart{E}{ste} proyecto logró construir un sistema funcional
de teleoperación donde un operador puede mover un brazo robótico
con la mano y ver cómo otro brazo, en otro lugar, replica exactamente
ese movimiento en tiempo real. Para lograrlo, se fabricaron sensores
de presión caseros con esponja y grafito, se programaron dos
microcontroladores ESP32 y se conectó todo a través de \ros{} Jazzy.
 
Uno de los resultados más satisfactorios fue la detección de
contacto: cuando el brazo esclavo choca contra un obstáculo, ambos
brazos se detienen automáticamente en menos de 100\,ms, protegiendo
tanto al robot como al entorno. Al retirar el obstáculo, el sistema
reanuda la operación por sí solo.
 
Los sensores artesanales resultaron sorprendentemente útiles para
detectar si hay contacto o no, aunque requieren calibración
individual y se descalibran con el tiempo y la temperatura. El
software de control demostró que es posible mover un brazo en
cualquier dirección cartesiana sin necesidad de programar
manualmente la geometría del robot, simplemente calculando el
Jacobiano desde la información de posición que ya publica \ros{}.
 
Los errores de seguimiento medidos (0{,}33--1{,}10°\,RMS) y el
retardo de $\approx$\,100--150\,ms en el robot físico son resultados
razonables para un sistema construido con hardware de bajo costo y
sin sensores de fuerza comerciales.
El código, firmware y datos están disponibles
públicamente~\cite{repo}.

\section*{Agradecimientos}
Los autores agradecen al Computational Robotics Lab del Tecnológico
de Monterrey y al Prof.\ Alberto Muñoz por el equipo robótico y
la orientación académica.

%% ── Bibliografía ─────────────────────────────────────────────
\begin{thebibliography}{11}
\bibitem{xarm_github}
  UFACTORY, ``xArm-Python-SDK,'' GitHub, 2024.
  [En línea]: \url{https://github.com/xArm-developer/xArm-Python-SDK}
\bibitem{siciliano2009}
  B.~Siciliano \textit{et al.},
  \textit{Robotics: Modelling, Planning and Control}. Springer, 2009.
\bibitem{spong2006}
  M.\,W.~Spong, S.~Hutchinson, M.~Vidyasagar,
  \textit{Robot Modeling and Control}. Wiley, 2006.
\bibitem{hogan1985}
  N.~Hogan, ``Impedance Control,''
  \textit{ASME J.\ Dyn.\ Sys.}, vol.\,107, 1985.
\bibitem{colgate1993}
  J.\,E.~Colgate, J.\,M.~Brown, ``Factors affecting the Z-width,''
  \textit{Proc.\ IEEE ICRA}, 1994.
\bibitem{murray1994}
  R.\,M.~Murray, Z.~Li, S.\,S.~Sastry,
  \textit{A Mathematical Introduction to Robotic Manipulation}.
  CRC Press, 1994.
\bibitem{niemeyer2008}
  G.~Niemeyer, C.~Preusche, G.~Hirzinger, ``Telerobotics,''
  \textit{Springer Handbook of Robotics}, 2008.
\bibitem{thompson2001}
  S.~Thompson \textit{et al.},
  ``Using a stationary camera\ldots,''
  \textit{IEEE Trans.\ SMC-A}, vol.\,31(1), pp.\,74--82, 2001.
\bibitem{safeea2023}
  M.~Safeea, P.~Neto,
  ``Minimum time trajectory planning for cobots
  using jerk-limited motion profiles,''
  \textit{IEEE Trans.\ Ind.\ Electron.}, vol.\,70, no.\,9,
  pp.\,9112--9121, Sep.\ 2023.
\bibitem{repo}
  ``Teleoperaci\'on con Sensado de Fuerza --- \xarm{},''
  GitHub, 2026.
  [En l\'inea]: \url{https://github.com/kurttco/TeleOp-Haptica.git}
\end{thebibliography}

\end{document}
%% ============================================================